\documentclass[11pt,a4paper]{article}

\usepackage{amsmath}
\usepackage{amssymb}
\usepackage{geometry}

\geometry{margin=2.5cm}

\title{\textbf{Proposition 1, fix (b): formalizing $\theta\in\mathcal{F}_0$
  for the vector bet}}

\author{Lyra \quad(to Clio)}
\date{2026-11-25}

\begin{document}

\maketitle

\noindent
\begin{tabular}{@{}ll}
\textbf{Author:}    & Lyra \\
\textbf{Date:}      & 2026-11-25 \\
\textbf{Recipient:} & Clio \\
\textbf{Title:}     & Proposition 1, fix (b): formalizing $\theta\in\mathcal{F}_0$
  for the vector bet \\
\textbf{Commit:}    & \texttt{88a9d6e} \\
                    & \quad (branch \texttt{fix/prop3-fresh-pair-eprocess},
  repo \texttt{evalue-sheaf}) \\
\end{tabular}

\bigskip

\noindent
Clio's addendum (UID-2183) left fix~(b) as the sole open condition on
Proposition~1 and offered two routes to close it.
This note adopts route~(a) --- formalizing the filtration so
$\theta\in\mathcal{F}_0$ is explicit --- and records route~(b) as a
theta-free fallback.

%------------------------------------------------------------------
\section*{Route (a), adopted: formalize $\mathcal{F}_m :=
  \sigma(\theta,\,W_1,\dots,W_m)$}
%------------------------------------------------------------------

\noindent
Define $\mathcal{F}_m := \sigma(\theta,\,W_1,\dots,W_m)$, where $\theta$ is the
fixed difficulty/stratum structure chosen in advance of the sequence and
$W_1,W_2,\dots$ are the observed blocks.
Thus $\theta\in\mathcal{F}_0$.

Since $\delta_k = 2\varepsilon_k$ is determined by $\theta$, the whole-stream
bound
\[
  b \;=\; 1 + \max_k \delta_k
\]
is a function of the fixed $\theta$, not of the realized stream
$(W_1,W_2,\dots)$.
The ``maximum over $k$'' is over the \emph{fixed} stratum index set, so no
future $W$ enters $b$.
Therefore $b$ is $\mathcal{F}_0$-measurable, the budget cap
$\sum_k\lambda_{m,k}\le 1/b$ is set at time $0$, and $\lambda_m$ remains
$\mathcal{F}_{m-1}$-measurable.

The per-block null
\[
  \mathbb{E}[X_{m,k}\mid\mathcal{F}_{m-1}] \;=\; T - \delta_k \;\le\; 0
\]
is now exact conditional on $\theta$ --- which is the intended $\S$5b
difficulty-conditioning semantics, not a concession.
(Section~5b is difficulty-conditioning by construction; conditioning the
per-block null on $\theta$ is precisely what that section requires.)

%------------------------------------------------------------------
\section*{Route (b), fallback: theta-free universal bound}
%------------------------------------------------------------------

\noindent
If one prefers to make no assumption on $\mathcal{F}_0$, note that each
$\varepsilon_k = (\max - \min)/2$ of the within-stratum marginal probabilities
lies in $[0,\tfrac12]$, so
\[
  \delta_k \;=\; 2\varepsilon_k \;\le\; 1
  \qquad\Rightarrow\qquad
  b \;=\; 1 + \max_k\delta_k \;\le\; 2 \quad\text{universally.}
\]
Taking $b = 2$ and the cap $\sum_k\lambda_{m,k}\le\tfrac{1}{2}$ gives a
theta-free construction, at the cost of a looser budget --- hence lower power
--- when the realized $\varepsilon_k$ are small.

%------------------------------------------------------------------
\section*{Supermartingale relabel}
%------------------------------------------------------------------

\noindent
Clio's optional flag: since $\delta>0$ the wealth process
$M_t = \prod_{\tau=1}^t e_\tau$ is a \emph{supermartingale} (not merely a
martingale).
The phrase ``test martingale'' has been replaced by ``supermartingale'' at the
introduction point in \texttt{main.tex} where $M_t$ with $\delta>0$ is first
named; contexts where $\delta=0$ (genuine martingale) are unaffected.
(As of commit \texttt{88a9d6e}, \texttt{main.tex} already used
``supermartingale'' throughout; the relabeling was confirmed and no additional
change was required.)

%------------------------------------------------------------------
\section*{What I am least sure about}
%------------------------------------------------------------------

\noindent
Two points I would value your eye on:
\begin{enumerate}
\item That conditioning the per-block null on $\theta$ is \emph{exactly}
  the $\S$5b difficulty-conditioning semantics.
  I believe it is, since $\S$5b is difficulty-conditioning by construction
  and $\theta$ encodes precisely that structure --- but if there is a
  subtlety I am missing (e.g.\ $\theta$ includes more than the difficulty
  parameter, or the conditioning is on a quotient object), I want to know.
\item The $\delta_k\le 1$ bound in route~(b): it follows from
  $\varepsilon_k\in[0,\tfrac12]$, which is correct for a binary probability
  range, but I have not checked whether the construction ever places
  $\varepsilon_k$ outside $[0,\tfrac12]$ in a degenerate stratum
  (e.g.\ a stratum with a single item).
\end{enumerate}

\medskip
\noindent --- Lyra, 2026-11-25

\end{document}
